\uuid{gmZF}
\titre{Propriété d’Archimède}
\auteur{bodin ; hanani ; mzari}
\organisation{exo7}
\licence{CC BY-NC-SA 4.0}
\datecreate{2020-01-20}
\niveau{L1}
\module{Analyse}
\chapitre{Propriétés de R}
\sousChapitre{Propriétés des nombres réels}
\difficulte{3}
\theme{Propriétés de nombres réels, propriété d’Archimède}
\exo7qcmid{lille-1.reels.015}

\contenu{
\begin{question}
\qclassification{120.03}
\qtags{motcle=propriété d’Archimède}
Quelles propriétés découlent de la propriété d’Archimède des réels (c'est-à-dire $\Rr$ est archimédien) ?
\begin{answers}
    \bad{$\exists x>0 \quad \forall n \in\Nn \quad x>n$}

    \good{$\forall x>0 \quad \exists n \in\Nn \quad n>x$}

    \good{$\forall \epsilon >0 \quad \exists n \in \Nn \quad 0 < \frac 1n < \epsilon$}

    \good{$\forall x>0 \quad \forall y>0 \quad \exists n \in \Nn \quad nx >y$}
\end{answers}
\begin{explanations}
La définition de la propriété d'Archimède est $\forall x>0 \quad \exists n \in\Nn \quad n>x$. Cela implique les deux autres assertions vraies.
\end{explanations}
\end{question}
}
